\subsection{Neutrino Vertex Resolution}\label{sec:VERTEX_RESOLUTION}
The vertex resolution is defined from the distance between the reconstructed vertex and the true vertex position. Note that this analysis reconstructs particle interactions as point--like objects and therefore reconstructs a vertex that is shifted on average to the center of the true particle track. For the cases in which a coordinate cannot be determined (applicable to x- and y- coordinates), the reconstructed coordinate is set to the center of the active volume. The vertex error is then defined as the difference between the reconstructed and true coordinates.

Given that the distributions are not Gaussian, we can still attempt to fit a Gaussian function to the vertex error distributions by ignoring the tails to extract a vertex resolution from the FWHM of the fit as a figure of merit. This is done by fitting a Gaussian function to the core of the distribution.

\begin{figure}[ht!]
    \centering
    \includegraphics[width=.49\linewidth]{Images/LowE/vertex_smearing_hd_1x2x6_centralapa_marley_official_coordinate_logy_hist1d.png}
    \includegraphics[width=.49\linewidth]{Images/LowE/purity_match_resolution_hd_1x2x6_centralapa_marley_official_values_counts_label_coordinate_x_logy_fit.png}
    \caption[width=1\linewidth]{\gls{hd} vertex error distributions. (left) All coordinate error distributions. (right) X coordinate error divided in high- and low-purity matched flash samples with core distribution fit.\label{fig:SIGNAL_VERTEX_RESOLUTION}}
\end{figure}

The results of this fitting procedure as a function of true neutrino energy are presented in~\cref{fig:SIGNAL_VERTEX_RESOLUTION} for both the neutrino-only and neutrino-plus-background samples.

\begin{table}[ht!]
    \centering
    \centering
Vertex Resolution~(cm))\\
\begin{tabular}{lccc}
\toprule
Configuration & \multicolumn{3}{c}{Vertex} \\
 & X & Y & Z \\
\midrule
\emph{HD Signal} & 0.72 ± 0.01 & 0.55 ± 0.01 & 0.64 ± 0.01 \\
\emph{HD Lateral} & 0.65 ± 0.03 & 0.68 ± 0.01 & 0.65 ± 0.01 \\
\emph{HD Central} & 0.71 ± 0.02 & 0.65 ± 0.01 & 0.65 ± 0.01 \\
\emph{VD Signal} & 0.69 ± 0.01 & 0.69 ± 0.01 & 0.67 ± 0.01 \\
\emph{VD Top} & 0.58 ± 0.03 & 0.66 ± 0.01 & 0.68 ± 0.02 \\
\bottomrule
\end{tabular}
    \caption[width=1\linewidth]{Vertex resolution comparison between all geometry and background configurations. The values are extracted from the FWHM of the Gaussian fit to the core of the vertex error distribution.\label{tab:VERTEX_RESOLUTION_COMPARISON}}
\end{table}

\Cref{tab:VERTEX_RESOLUTION_COMPARISON} summarizes the vertex resolution for all geometry and background configurations. The vertex resolution is found to be between \SI{0.5}{\centi\meter} and \SI{0.75}{\centi\meter}. The fact that all coordinates and configurations result in a similar vertex resolution points to the fact that the vertex resolution is dominated by the intrinsic smearing of the point-like reconstruction approach, which places the reconstructed vertex at the center of the neutrino-induced particle track. As the neutrino energy increases, the average particle track size also increases, leading to a larger displacement between the true interaction vertex and the reconstructed vertex. Currently, there is no low--energy physics analysis in DUNE that would strongly benefit from a vertex resolution better than a few centimeters. For example, in the case of solar neutrinos, where external cavern--gammas are one of the major backgrounds, the Compton scattering length is on the order of tens of centimeters, which is much larger than the achieved vertex resolution. But, if needed, the vertex resolution could be improved by implementing a more sophisticated reconstruction algorithm that takes into account the topology of the event.

\begin{figure}[ht!]
    \centering
    \includegraphics[width=.49\linewidth]{Images/Reco/vertex_reconstruction_efficiency_hd_1x2x6_centralapa_marley_official_values_efficiency_energy_variable_x_sigma_true_tolerance_3_scan.png}
    \includegraphics[width=.49\linewidth]{Images/Reco/vertex_reconstruction_efficiency_vd_1x8x14_3view_30deg_nominal_marley_official_values_efficiency_energy_variable_x_sigma_true_tolerance_3_scan.png}
    \caption[width=1\linewidth]{\gls{hd} (left) and \gls{vd} (right) vertex reconstruction efficiency as a function of true neutrino energy for different vertex error tolerances.\label{fig:VERTEX_RECONSTRUCTION_EFFICIENCY}}
\end{figure}

\Cref{fig:VERTEX_RECONSTRUCTION_EFFICIENCY} shows the vertex reconstruction efficiency as a function of true drift coordinate for different neutrino energies. The efficiency is defined as the fraction of events for which the reconstructed vertex is within a certain distance (3 widths of the core error distribution) from the true vertex position. As expected, the efficiency increases with increasing neutrino energy, as higher energy events tend to produce more ionization charge and scintillation light, leading to better vertex reconstruction accuracy.

